\section{行列とベクトル}\label{節：行列とベクトル}
\subsection{値を一組にする}

二つの数が
\[
  x'=2x+y,\qquad y'=x+2y
\]
によって同時に変わるとしよう.
$(x,y)=(1,0)$からは$(2,1)$,\,次に$(5,4)$へ進む.
二つ目を計算するときに,\,更新済みの$x'$を代入してはいけない.
二つの式はどちらも更新前の$x,y$を使う.

この二つの値を縦に並べた
\[
  \boldsymbol v=\begin{pmatrix}x\\y\end{pmatrix}
\]
を\textbf{列ベクトル}と呼ぶ.
平面上の点$(x,y)$の位置を表す矢印としても読める.
ベクトルの各数を成分といい,\,和と定数倍は成分ごとに計算する.
\[
 \begin{pmatrix}x\\y\end{pmatrix}
 +\begin{pmatrix}u\\v\end{pmatrix}
 =\begin{pmatrix}x+u\\y+v\end{pmatrix},
 \qquad
 c\begin{pmatrix}x\\y\end{pmatrix}
 =\begin{pmatrix}cx\\cy\end{pmatrix}.
\]
等号も,\,対応する成分どうしの等しさを表す.
三つ以上の数を並べる場合にも同じように定める.

\subsection{行列は更新を記録する}

数を長方形に並べたものを\textbf{行列}という.
横の並びを行,\,縦の並びを列と呼ぶ.
例えば二行二列の行列
\[
 A=\begin{pmatrix}2&1\\1&2\end{pmatrix}
\]
を列ベクトルに掛ける積を
\[
 A\begin{pmatrix}x\\y\end{pmatrix}
 =\begin{pmatrix}2x+y\\x+2y\end{pmatrix}
\]
と定める.
第一行は更新後の第一成分を,\,第二行は第二成分を指定する.
冒頭の二つの更新を,\,一つの積で書けた.

一般には
\[
 \begin{pmatrix}p&q\\r&s\end{pmatrix}
 \begin{pmatrix}x\\y\end{pmatrix}
 =\begin{pmatrix}px+qy\\rx+sy\end{pmatrix}.
\]
また,\,各成分を展開すると
\[
 A(c\boldsymbol u+d\boldsymbol v)
 =cA\boldsymbol u+dA\boldsymbol v
\]
となる.
先に足してから更新しても,\,別々に更新してから足しても同じである.
この性質を\textbf{線形性}という.
後で初期状態をいくつかの部分に分けて追う際の根拠になる.

\subsection{積と累乗}

同じ行列で二回更新することを$A^2$で表す.
例えば
\[
 A^2\begin{pmatrix}x\\y\end{pmatrix}
 =A\begin{pmatrix}2x+y\\x+2y\end{pmatrix}
 =\begin{pmatrix}5x+4y\\4x+5y\end{pmatrix}
\]
だから,\,先ほどの$A$では
\[
 A^2=\begin{pmatrix}5&4\\4&5\end{pmatrix}
\]
である. 各成分の二乗ではない.

まず$B$,\,次に$A$で更新する操作を$AB$と書く.
積の右側から先に作用する.
一般に$AB=BA$とは限らない.
例えば横方向に2倍する
$A=\begin{pmatrix}2&0\\0&1\end{pmatrix}$と,
$B=\begin{pmatrix}1&1\\0&1\end{pmatrix}$では
\[
 AB=\begin{pmatrix}2&2\\0&1\end{pmatrix},
 \qquad BA=\begin{pmatrix}2&1\\0&1\end{pmatrix}.
\]
順序を変えると,\,第二成分が第一成分に寄与する倍率が変わる.

何も変えない更新
\[
 I=\begin{pmatrix}1&0\\0&1\end{pmatrix}
\]
を\textbf{単位行列}といい,\,更新を一度もしないことを$A^0=I$と定める.
従って,\,$\boldsymbol v_{n+1}=A\boldsymbol v_n$なら
\[
 \boldsymbol v_n=A^{n-1}\boldsymbol v_1
\]
である. あとは,\,この累乗をどう読むかが問題になる.

\subsection{座標と基底}\label{節：座標と基底}

平面の点を表す二つの数は,\,選んだ基準方向によって変わる.
例えば
\[
 \boldsymbol u=\begin{pmatrix}1\\1\end{pmatrix},
 \qquad
 \boldsymbol w=\begin{pmatrix}1\\-1\end{pmatrix}
\]
を基準にすれば,
\[
 \begin{pmatrix}x\\y\end{pmatrix}
 =\frac{x+y}{2}\boldsymbol u
  +\frac{x-y}{2}\boldsymbol w.
\]
点は同じでも,\,新しい座標は$\frac{x+y}{2},\,\frac{x-y}{2}$となる.
\NewFigBasis

ベクトルを定数倍して足したものを\textbf{線形結合}という.
その線形結合が零になるのが,\,すべての係数が零の場合だけであるとき,\,元のベクトルの組を\textbf{一次独立}という.
例えば二本の場合は,\,零でなく互いに比例しないことに対応する.
数列でも,\,各項で定数倍と和を取ることで,\,同じ言葉を使う.

一般に,\,どんなベクトルも一意に線形結合で表せる基準ベクトルの組を\textbf{基底}という.
二成分の実ベクトルでは,\,零でなく互いに比例しない二本を選べばよい.
実際,\,$\boldsymbol u=\begin{pmatrix}u_1\\u_2\end{pmatrix}$,
$\boldsymbol w=\begin{pmatrix}w_1\\w_2\end{pmatrix}$に対して
\[
 c\boldsymbol u+d\boldsymbol w
 =\begin{pmatrix}x\\y\end{pmatrix}
\]
を成分に分けると,\,$c,d$の連立一次方程式になる.
係数の組が比例しなければ一意に解ける.

二行二列の行列
$P=\begin{pmatrix}u_1&w_1\\u_2&w_2\end{pmatrix}$
では,\,列が基底になる条件は
\[
 \det P=u_1w_2-w_1u_2\ne0
\]
である. この数を\textbf{行列式}という.
この場合,\,更新を逆に戻す行列
\[
 P^{-1}
 =\frac1{\det P}
 \begin{pmatrix}w_2&-w_1\\-u_2&u_1\end{pmatrix}
\]
があり,\,$P^{-1}P=PP^{-1}=I$を満たす.
$P^{-1}$を\textbf{逆行列}という.
元の座標を$\boldsymbol v$,\,新しい座標を$\boldsymbol c$とすれば
\[
 \boldsymbol v=P\boldsymbol c,\qquad
 \boldsymbol c=P^{-1}\boldsymbol v.
\]
座標を変える操作と元へ戻す操作を,\,行列で書けた.

\begin{bookpoint}
行列の累乗を計算する前に,\,\textbf{更新が簡単に見える基底を選べないか}を考える.
座標軸が直角である必要も,\,基準ベクトルの長さが1である必要もない.
\end{bookpoint}
